\begin{lemma}
Let $H$ be a finite $k$-uniform hypergraph and let $K,N$ be natural
numbers. Suppose $C\subseteq E(H)$ has a proper coloring with $K$ colors
and its canonical residual on $E(H)\setminus C$ has maximum degree at
most $N$. Then $H$ has a proper edge-coloring with $K+kN+1$ colors.
\end{lemma}
\begin{proof}
The canonical residual is $k$-uniform because its edge sets are unchanged.
By \noderef{UniformHypergraphLineDegreeBound}, each degree in its line
graph is at most $k(N-1)\le kN$, with natural subtraction understood.
By \noderef{FiniteGraphGreedyColoring}, it has a proper vertex-coloring
with the positive number $kN+1$ of colors. Use these on residual edges,
adding $K$ to their labels, and retain the original labels on $C$.
Intersecting edges in $C$ have distinct colors by assumption;
intersecting residual edges are adjacent in the line graph and thus
receive distinct colors. A pair with one edge in each part has distinct
colors because its labels lie on opposite sides of $K$. This is the
required \noderef{ChromaticIndexAtMost} coloring.
\end{proof}
